\chapter{SDR та RF-передній край: архітектурна математика}
\label{ch:sdr-rf-architecture}

\sourcebox{Основні джерела для цієї доріжки: переклад translation system для Akeela--Dezfouli, \emph{Software-defined Radios: Architecture, State-of-the-art, and Challenges}, arXiv:1804.06564; раніше staged PySDR Ukrainian RST; базовий DSP-розділ цього довідника. Цей розділ є інтегрованою навчальною нормалізацією, а не дослівним перекладом survey.}

\section{Сигнальний шлях SDR}

Програмно-визначене радіо (SDR) треба читати не як окрему плату або конкретний набір інструментів, а як архітектурний принцип: частину радіосистеми, яку раніше фіксували апаратно, переносять у цифрову обробку сигналів. Типовий шлях сигналу з боку приймання:

\begin{center}
антена $\to$ RF-передній край $\to$ змішувач/локальний осцилятор $\to$ IF або базова смуга $\to$ ADC $\to$ цифровий передній край $\to$ DSP/FPGA/GPU/GPP.
\end{center}

У зворотному напрямку для передавання з'являються DAC, цифрове перетворення догори, підсилення і фільтрація. Математично найзручніше працювати з комплексною базовою смугою. Якщо $x_{BB}(t)=I(t)+jQ(t)$, то дійсний RF-сигнал з несною $f_c$ можна записати як
\begin{equation}
  x_{RF}(t)=\Re\{x_{BB}(t)e^{j2\pi f_c t}\}
  = I(t)\cos(2\pi f_c t)-Q(t)\sin(2\pi f_c t).
\end{equation}
Ця формула є ключем до I/Q-представлення: дві дійсні послідовності $I[n]$ і $Q[n]$ описують амплітуду та фазу вузькосмугового сигналу без явного моделювання швидкої несної.

\section{Дискретизація і цифрове перетворення донизу}

Нехай неперервний сигнал $x(t)$ дискретизується з частотою $f_s=1/T_s$:
\begin{equation}
  x[n] = x(nT_s).
\end{equation}
Для смугообмеженого сигналу потрібна частота дискретизації, достатня для уникнення аліасингу після антиаліасингової фільтрації. У реальній системі важливо відрізняти три твердження:

\begin{enumerate}[label=\alph*)]
\item теорема Найквіста-Шеннона описує ідеальний математичний випадок;
\item реальний ADC має скінченну розрядність, джитер, нелінійність і обмежений SFDR;
\item decimation допустиме тільки після низькочастотної фільтрації, яка пригнічує спектральні копії.
\end{enumerate}

Цифровий перетворювач донизу (DDC) зазвичай виконує множення на цифрову комплексну експоненту, фільтрацію і зменшення частоти дискретизації:
\begin{align}
  y[n] &= x[n]e^{-j2\pi f_0 n/f_s},\\
  z[m] &= \sum_{k} h[k] y[mM-k],
\end{align}
де $M$ - коефіцієнт децимації, а $h[k]$ - імпульсна характеристика низькочастотного фільтра.

\section{ADC/DAC як математичні обмеження}

Ідеалізований $N$-бітовий рівномірний квантувач має крок
\begin{equation}
  \Delta = \frac{V_{\max}-V_{\min}}{2^N}.
\end{equation}
Якщо похибка квантування моделюється рівномірним шумом на $[-\Delta/2,\Delta/2]$, то її дисперсія дорівнює
\begin{equation}
  \sigma_q^2 = \frac{\Delta^2}{12}.
\end{equation}
Для ідеального синусоїдального сигналу часто використовують оцінку
\begin{equation}
  \mathrm{SNR}_{\mathrm{dB}} \approx 6.02N + 1.76.
\end{equation}
Це не паспортний опис конкретного пристрою, а sanity check: якщо розрахункові вимоги до динамічного діапазону набагато перевищують можливості ADC, проблема не розв'язується перекладом або програмуванням.

\section{Вибір обчислювальної платформи як задача компромісу}

У SDR-архітектурі одна й та сама математична операція може бути реалізована на різних рівнях. Наприклад, FIR-фільтр
\begin{equation}
  y[n] = \sum_{k=0}^{L-1} b_k x[n-k]
\end{equation}
можна виконувати на CPU, GPU, DSP або FPGA. Вибір залежить не від назви платформи, а від латентності, пропускної здатності, енергії, гнучкості та складності верифікації.

\begin{longtable}{L{0.17\linewidth}L{0.25\linewidth}L{0.25\linewidth}L{0.22\linewidth}}
\toprule
\textbf{Платформа} & \textbf{Сильна сторона} & \textbf{Типове обмеження} & \textbf{Коли розглядати першою}\\
\midrule
GPP/CPU & простота розробки, Python/C++ екосистема & латентність і пропускна здатність & прототип, контрольні алгоритми, offline аналіз.\\
GPU & висока паралельність для великих блоків & передавання даних і latency overhead & батчова обробка, ML/CV, великі FFT.\\
DSP & спеціалізовані MAC-операції & менша гнучкість екосистеми & потокова DSP-обробка з помірною латентністю.\\
FPGA & детермінована latency, pipeline, паралельність & складність розробки і верифікації & DDC/DUC, фільтри, синхронізація, високошвидкісні тракти.\\
ASIC & ефективність, площа, потужність & майже нульова гнучкість після виготовлення & масове виробництво стабільного дизайну.\\
\bottomrule
\end{longtable}

\section{Що перевіряти у перекладі SDR-матеріалу}

Для SDR-текстів найбільш руйнівні не стилістичні помилки, а зсуви одиниць і частотних позначень. Мінімальний QA-набір:

\begin{enumerate}[label=\arabic*.]
\item \textbf{Hz, rad/s, cycles/sample.} Перевірити, чи $2\pi$ не зник і не з'явився зайвий раз.
\item \textbf{Baseband vs RF.} Не перекладати baseband як ``основна смуга'' в одному місці і ``базова смуга'' в іншому без рішення словника.
\item \textbf{Sample, sample rate, sampling.} У DSP sample = відлік; sampling = дискретизація; sample rate = частота дискретизації.
\item \textbf{Decimation/interpolation.} Децимація - це не просто викидання відліків; перед нею потрібна фільтрація.
\item \textbf{Power vs amplitude.} Для потужності $10\log_{10}$, для амплітуди $20\log_{10}$.
\end{enumerate}

Цей розділ достатній як міст між базовим Fourier/DSP-ядром і більш спеціальними SDR-перекладами. Повний survey треба тримати окремо, бо його цінність - у термінології та архітектурній карті, а не у dated переліках апаратних платформ.
